% Extracted from AdaptiveAI(5).tex; manuscript snapshot, not preregistration.
\subsection{Evidence Extraction and Methodological Appraisal}

The extraction framework preserves the experimental
context required to interpret reported findings.
Six information categories are distinguished:

\begin{enumerate}
	\item \textbf{Study characteristics (EE1):}
	publication identifier, year, research objective,
	application domain, and learning paradigm.
	
	\item \textbf{Adaptation mechanism (EE2):}
	update strategy, adaptation trigger, learning
	objective, supervision, historical-data access,
	and preservation mechanisms.
	
	\item \textbf{Environmental change and evaluation
		protocol (EE3):} distribution shifts, task
	transitions, temporal ordering, task-boundary
	information, and naturally occurring or
	experimentally constructed changes.
	
	\item \textbf{Comparator characteristics (EE4):}
	adaptation status, architecture, initialisation,
	pre-training, information access, and
	computational budgets.
	
	\item \textbf{Experimental outcomes (EE5):}
	performance, capability-preservation, and
	resource measurements in their original
	definitions and units.
	
	\item \textbf{Evidence provenance (EE6):}
	source publication, experimental configuration,
	table, figure, equation, textual passage, and
	PDF page number where available.
\end{enumerate}

Multiple adaptation categories may be assigned to
a single system. Machine-assisted identification
of candidate evidence is distinguished from
source-level numerical verification; automated
indicators do not establish the presence or absence
of valid experimental evidence.

The proposed methodological appraisal comprises
seven claim-specific domains. These reflect concerns
arising from experimental design, computational
reporting, benchmark variability, and reproducibility
in machine learning
\cite{wohlin2012experimentation,bouthillier2021variance,
	pineau2021improving,henderson2020systematic}.

\begin{description}
	\item[\textbf{Q1: Environmental-change validity.}]
	Relevance of the experimental change scenario,
	task transitions, temporal ordering, and
	information about environmental changes.
	
	\item[\textbf{Q2: Comparator fairness.}]
	Comparability of architecture, initialisation,
	pre-training, information access, adaptation
	budgets, and evaluation conditions.
	
	\item[\textbf{Q3: Experimental validity.}]
	Separation of training and evaluation
	information, avoidance of data leakage and
	future-information access, and consistency
	of experimental conditions.
	
	\item[\textbf{Q4: Capability-preservation assessment.}]
	Explicit post-adaptation evaluation of previously
	acquired capabilities using interpretable
	retention measurements.
	
	\item[\textbf{Q5: Resource-accounting validity.}]
	Completeness and comparability of hardware,
	execution-time, memory, parameter, energy,
	and information-access measurements.
	
	\item[\textbf{Q6: Statistical reliability.}]
	Adequacy of repeated trials, random seeds,
	task-order sensitivity, uncertainty estimates,
	and treatment of dependent observations.
	
	\item[\textbf{Q7: Reproducibility.}]
	Availability of implementation details, data,
	code, hyperparameters, and evaluation procedures
	necessary to understand or reproduce the
	reported experiments.
\end{description}

Each applicable appraisal domain is classified as
\emph{low concern}, \emph{high concern},
\emph{unclear concern}, or \emph{not applicable},
according to methodological adequacy, identified
limitations, information sufficiency, and relevance
to the claim.

Judgements remain claim-specific rather than being
aggregated into a composite quality score. The
appraisal framework is review-specific and has not
been independently validated as a risk-of-bias
instrument. Aggregate appraisal results are not
reported because formal domain-level assessments
remain incomplete.

Evidence availability is distinguished from
methodological validity. Unresolved and
non-applicable assessments are distinguished from
verified absence of reporting. Missing retention
or resource measurements do not imply zero
forgetting or zero expenditure.
